\section*{ЗАКЛЮЧЕНИЕ}
\phantomsection
\addcontentsline{toc}{section}{Заключение}

В ходе выполнения курсовой работы проведён численный анализ влияния параболического дискретного микрорельефа на рабочие характеристики радиального подшипника скольжения. В процессе работы решены следующие задачи:

\begin{enumerate}
    \item Выполнен обзор теоретических основ гидродинамической теории смазки и современных работ, посвящённых лазерному текстурированию поверхностей трения. Показано, что основным механизмом положительного действия микрорельефа является микрогидродинамический эффект, возникающий на кромках углублений за счёт локального сужения и расширения смазочного зазора.

    \item Реализована методика численного решения уравнения Рейнольдса для подшипника конечной длины с применением метода конечных разностей и итерационной процедуры Гаусса--Зейделя с верхней релаксацией. Моделирование выполнено на сетке $500 \times 500$ узлов с критерием сходимости $\delta < 10^{-5}$.

    \item Рассчитаны поля давления и интегральные характеристики (нагрузочная способность $F$, коэффициент трения $\mu$, расход смазки $Q$) для гладкого подшипника и подшипника с параболическими углублениями (тип~2) в диапазоне эксцентриситетов $\varepsilon = 0{,}05 - 0{,}80$. Геометрия подшипника: $R = 50$~мм, $c = 70$~мкм, $L = 80$~мм.

    \item Проведено сравнение гладкого и текстурированного подшипников. При характерном рабочем эксцентриситете $\varepsilon = 0{,}6$ получены следующие результаты:
    \begin{itemize}
        \item нагрузочная способность возросла на $3{,}18$\%: с $F = 10685{,}6$~Н до $F = 11025{,}8$~Н;
        \item коэффициент трения уменьшился на $3{,}80$\%: с $\mu = 0{,}01507$ до $\mu = 0{,}01450$;
        \item расход смазки увеличился на $0{,}42$\%: с $Q = 0{,}6739$~л/с до $Q = 0{,}6768$~л/с.
    \end{itemize}
\end{enumerate}

Полученные данные свидетельствуют о том, что нанесение параболических углублений с параметрами $a = 2{,}5$~мм, $b = 2{,}3$~мм, $h_p = 10$~мкм ($h_p/c = 0{,}1$) приводит к улучшению рабочих характеристик подшипника. Одновременное повышение несущей способности и снижение коэффициента трения при незначительном росте расхода смазки указывает на эффективность данного типа микрорельефа.

Параболический профиль углублений обеспечивает квадратичное изменение глубины от центра к периферии, что формирует более выраженные градиенты давления на кромках по сравнению с эллипсоидальным профилем, однако исключает резкие скачки, свойственные цилиндрическим углублениям. Такой характер изменения глубины создаёт устойчивый дополнительный вклад в несущую способность, что делает параболический профиль перспективным решением для подшипников, работающих в стационарных режимах.

На основании результатов расчёта можно рекомендовать применение параболического микрорельефа для подшипников скольжения нефтегазового оборудования, эксплуатируемого в условиях жидкостного трения. Наибольший эффект достигается при высоких эксцентриситетах ($\varepsilon > 0{,}5$), где прирост нагрузочной способности и снижение трения выражены в наибольшей степени.
